% ৪.৫ ডিজাইন: স্ট্রাকচার ও লেআউট
\section{ডিজাইন: স্ট্রাকচার ও লেআউট}

\begin{learningbox}
এই টপিক শেষে শিক্ষার্থী পারবে:
\begin{itemize}
\item layout, wireframe ও responsive design-এর ধারণা ব্যাখ্যা করতে।
\item fixed, fluid ও responsive layout আলাদা করতে।
\item ভালো web layout-এর basic design principle লিখতে।
\end{itemize}
\end{learningbox}

\subsection{লেআউটের ধারণা}
ওয়েব পেজে header, navigation, content, sidebar, image, footer ইত্যাদি অংশ কোথায় থাকবে, সেটির পরিকল্পনাকে layout বলা হয়। ভালো layout page-কে পরিষ্কার, পড়ার উপযোগী ও ব্যবহারকারী-বান্ধব করে।

\begin{center}
\fbox{\parbox{0.9\textwidth}{
\centering
\textbf{Page layout wireframe placeholder}\\[4pt]
Header\\
Navigation menu\\
Main content \hspace{1cm} Sidebar\\
Footer
}}
\end{center}

\subsection{লেআউটের ধরন}
\begin{center}
\begin{tabular}{|p{0.22\textwidth}|p{0.34\textwidth}|p{0.30\textwidth}|}
\hline
\textbf{Layout} & \textbf{ধারণা} & \textbf{সীমাবদ্ধতা/সুবিধা} \\
\hline
Fixed & width fixed থাকে & desktop-এ predictable, mobile-এ সমস্যা \\
\hline
Fluid & width percentage ভিত্তিক & screen অনুযায়ী ছড়ায়, control কম \\
\hline
Responsive & screen size অনুযায়ী layout বদলায় & modern website-এর জন্য উপযোগী \\
\hline
\end{tabular}
\end{center}

\begin{keypointbox}{Design principle}
\begin{itemize}
\item Contrast: গুরুত্বপূর্ণ অংশ আলাদা করে দেখানো।
\item Alignment: text ও element সুশৃঙ্খল রাখা।
\item Repetition: একই style বারবার ব্যবহার করে consistency রাখা।
\item Proximity: সম্পর্কিত বিষয় কাছাকাছি রাখা।
\end{itemize}
\end{keypointbox}

\begin{examplebox}{বাস্তব প্রয়োগ}
শিক্ষাপ্রতিষ্ঠানের website-এ top navigation-এ Home, Notice, Routine, Result, Contact রাখা যায়। Main area-তে latest notice এবং sidebar-এ important link রাখলে ব্যবহারকারী দ্রুত তথ্য পায়।
\end{examplebox}

\begin{practicebox}
\begin{enumerate}
\item Responsive design কী?
\item Wireframe কেন দরকার?
\item Fixed ও responsive layout-এর পার্থক্য লেখ।
\end{enumerate}
\end{practicebox}
